\documentclass[11pt,a4paper]{article}
\usepackage[T1]{fontenc}
\usepackage{lmodern}
\usepackage[margin=26mm]{geometry}
\usepackage{amsmath,amssymb,amsthm,booktabs,array,longtable}
\usepackage{microtype,multicol}
\usepackage[hidelinks]{hyperref}
\setlength{\parskip}{3pt}
\setlength{\emergencystretch}{2em}
\raggedbottom
\allowdisplaybreaks[1]
\clubpenalty=10000
\widowpenalty=10000
\newtheorem{theorem}{Theorem}[section]
\newtheorem{lemma}[theorem]{Lemma}
\newtheorem{corollary}[theorem]{Corollary}
\theoremstyle{remark}
\newtheorem{remark}[theorem]{Remark}
\newcommand{\Z}{\mathbb Z}
\newcommand{\Q}{\mathbb Q}
\newcommand{\N}{\mathbb N}
\newcommand{\F}{\mathbb F}
\newcommand{\ps}[1]{#1[[x]]}
\newcommand{\Rev}{\operatorname{Rev}}
\newcommand{\rev}{^{\langle-1\rangle}}
\newcommand{\Plus}{\operatorname{Plus}}
\newcommand{\Minus}{\operatorname{Minus}}
\newcommand{\coeff}[2]{[x^{#1}]#2}
\newcommand{\E}{\mathcal E}
\Urlmuskip=0mu plus 1mu
\hypersetup{pdftitle={A378580: Square-index residues for theta quotient substitution},pdfauthor={Dev Goyal},pdfsubject={Separate manuscript with AI disclosure and Lean proof map}}
\title{A378580\\[4pt]\large Square-index residues for theta quotient substitution}
\author{Dev Goyal}
\date{24 October 2026}
\begin{document}
\maketitle
\begin{center}
\fbox{\begin{minipage}{0.9\linewidth}\small
\textbf{AI disclosure.} ChatGPT (OpenAI) generated the mathematical arguments,
Lean proofs, and this manuscript in response to Dev Goyal's requests.
Dev Goyal is the sole human author. The source sequence and conjecture or
observation are credited to Paul D. Hanna.
\textbf{Independent human review is pending.}
\end{minipage}}
\end{center}
\begin{abstract}
Let $\Theta=1+2\sum_{j\ge1}x^{j^2}$. The unique constant-one integral solution of $A(x/A)=\Theta$ satisfies $A\equiv\Theta\pmod4$. We construct the literal unit quotient and prove that the positive-index residue is two precisely at squares and zero elsewhere.
\end{abstract}

\section{Source, notation, and scope}

The source is Paul D. Hanna's OEIS entry A378580 \cite{oeis580},
retrieved on 24 October 2026. Its initial sequence index is 0.
The constant coefficient is one; the residue classification below
is restricted to positive indices.

All series are formal. The set $\N$ includes zero, and $[x^n]U$ denotes a
coefficient. A normalized series is $x+O(x^2)$; $O(x^r)$ means divisibility
by $x^r$, over any coefficient ring used below. We write $U\circ V=U(V(x))$
only when $V(0)=0$. Compositional inversion, denoted $\Rev(U)$ or $U\rev$,
is distinct from multiplicative inversion of a unit. Congruences are
coefficientwise; ordinary products and compositional iterates are distinct.

In this manuscript $\Theta(x)=1+2\sum_{j\ge1}x^{j^2}$. The chosen source is
$A(x/A)=\Theta$. The auxiliary construction treats both signs.

\section{Formal algebra used in the proof}

\begin{lemma}[Coefficientwise construction]\label{lem:contraction}
Let $K$ be a commutative ring. Suppose an operator $\Phi$ preserves
normalized series and, for $r\ge2$, satisfies
\[
 G-H=O(x^r)\quad\Longrightarrow\quad
 \Phi(G)-\Phi(H)=O(x^{r+1}).
\]
Then $\Phi$ has exactly one normalized fixed point in $\ps K$.
More generally, a residual equation whose degree-$r$ condition determines
the next unknown coefficient with multiplier one has a unique solution
once its initial coefficients are fixed and its lower-degree conditions
hold identically.
\end{lemma}

\begin{proof}
Begin with $G_0=x$ and set $G_{m+1}=\Phi(G_m)$. Normalization gives
$G_1-G_0=O(x^2)$. Inductively,
$G_{m+1}-G_m=O(x^{m+2})$, so each coefficient eventually stabilizes.
Define $G$ using those stable coefficients. The displayed implication
ensures that each coefficient of $\Phi(G)$ equals the corresponding stable
coefficient, hence $G=\Phi(G)$. If two fixed points agree below degree
$r$, they agree through degree $r$; induction proves uniqueness.

For the residual version, set the next coefficient to the negative of the
known residual term. Its multiplier is one, so the construction uses only
ring operations. No earlier equation changes, and every equation is
eventually imposed. The same coefficient comparison proves uniqueness.
\end{proof}

\begin{lemma}[Units and reversion]\label{lem:inverse}
Every constant-one series has a unique constant-one multiplicative
inverse. Every normalized series has a unique normalized two-sided
compositional inverse. These inversions, formal differentiation, and
valid substitution commute with coefficient ring homomorphisms.
\end{lemma}

\begin{proof}
For $UV=1$ with $U(0)=1$, the coefficient of degree $n>0$ is $v_n$ plus
a finite expression in earlier coefficients of $V$, which determines
$v_n$ without division. For normalized $G$, the coefficient of degree
$n>1$ in $G(H)$ is $h_n$ plus an expression in earlier coefficients of
$H$, giving a unique normalized right inverse $H$.

To see the second inverse identity, construct a normalized right inverse
$J$ of $H$. Associativity gives
$G=G\circ(H\circ J)=(G\circ H)\circ J=J$, and hence $H\circ G=x$.
The same association proves uniqueness. Each coefficient of a valid
composition is a finite polynomial in coefficients. Differentiation
multiplies the degree-$(n+1)$ coefficient by $n+1$. Ring homomorphisms
therefore preserve these operations and equations; uniqueness identifies
the mapped inverses with the inverses of the mapped series.
\end{proof}

\begin{lemma}[Doubling and substitution]\label{lem:double}
For $U,V\in(\Z/4\Z)[[x]]$,
\[
 2U=2V\quad\Longleftrightarrow\quad
 U\equiv V\pmod2.
\]
If $B,D$ have zero constant coefficient and $B\equiv D\pmod2$, then
$2H(B)=2H(D)$ for every series $H$ over $\Z/4\Z$.
\end{lemma}

\begin{proof}
For residues $u,v$ modulo four, $2u=2v$ exactly when their reductions modulo
two agree. Apply this fact to each coefficient. Reduction modulo two
commutes with substitution, giving the second assertion. In particular,
this lemma never cancels two inside $\Z/4\Z$.
\end{proof}

\section{Construction and square support for A378580}

\subsection{A general integral construction}

To treat the selected source, we construct a pair of auxiliary substitution
equations. For a constant-one target, these equations are
\begin{equation}\label{eq:thetaoriginal}
 A_-(x/A_-(x))=\Theta(x),\qquad
 A_+(xA_+(x))=\Theta(x),\qquad A_\pm(0)=1.
\end{equation}
We first solve them with an arbitrary constant-one target $T$ in place
of $\Theta$, over any commutative coefficient ring $K$.

\begin{lemma}\label{lem:normalroot}
For each $T\in\ps K$ with $T(0)=1$, there is exactly one normalized $G$
satisfying
\begin{equation}\label{eq:normalroot}
 G\circ G=GT.
\end{equation}
\end{lemma}

\begin{proof}
Let $E(G)=G\circ G-GT$. Its coefficients in degrees zero and one vanish
for every normalized $G$. Suppose $G-H=d x^r+O(x^{r+1})$, with $r\ge2$.
Decompose the difference of self-compositions as
\[
 G\circ G-H\circ H=(G-H)\circ G+(H\circ G-H\circ H).
\]
Because $G=x+O(x^2)$, the first summand is
$d x^r+O(x^{r+1})$. Write $H=x+x^2W$. The second summand is
$(G-H)+G^2W(G)-H^2W(H)$. The difference after $G-H$ has order at least
$r+1$: factor $G^2-H^2$, and use that substitution in any series
preserves an agreement modulo $x^r$. Thus the second summand also is
$d x^r+O(x^{r+1})$.

Finally, $(G-H)T=d x^r+O(x^{r+1})$, since $T(0)=1$. Therefore
\[
 E(G)-E(H)=d x^r+O(x^{r+1}).
\]
The degree-$r$ residual has multiplier $2-1=1$ on the next unknown
coefficient. Lemma~\ref{lem:contraction}'s triangular construction proves
existence and uniqueness with no division, over the stated ring.
\end{proof}

Write this unique $G$ as $G=xU$. Then $U(0)=1$ and
\[
 G\circ G=G\,U(G)=GT.
\]
Cancelling $G$ is valid even when $K$ has zero divisors: multiplication by
$x$ is injective coefficientwise, and $U$ is a unit. It follows that
$U(xU)=T$. Conversely, any constant-one solution of $U(xU)=T$ gives a
normalized solution $xU$ of \eqref{eq:normalroot}. Consequently there
exists exactly one such $U$; denote it by $\Plus(T)$.

Define
\[
 \Minus(T)=\Plus(T^{-1})^{-1},
\]
where both occurrences of inverse here are multiplicative unit inverses.
If $U=\Plus(T^{-1})$, then $U(xU)=T^{-1}$.
Substitution respects products, so $U^{-1}(xU)=T$. Since
$x/\Minus(T)=xU$, this is
\[
 \Minus(T)\bigl(x/\Minus(T)\bigr)=T.
\]
Conversely, if constant-one $V$ satisfies $V(x/V)=T$, then inversion of
that identity gives $V^{-1}(xV^{-1})=T^{-1}$. The uniqueness of
$\Plus(T^{-1})$ identifies $V^{-1}=U$ and hence $V=\Minus(T)$.
We have proved existence and uniqueness for both source equations
over every commutative ring, including $\Z$.

The quotient is literal: $V(x)\,(xV(x)^{-1})=x$. It is not interpreted as a
compositional inverse. Coefficient maps preserve the two source equations
and their unit inverses, so uniqueness also proves that $\Plus$ and
$\Minus$ commute with reduction of coefficients.

\subsection{The common square-support congruence}

\begin{theorem}\label{thm:theta}
Let $A_-=\Minus(\Theta)$ and $A_+=\Plus(\Theta)$ over $\Z$.
These are the unique constant-one integral solutions of
\eqref{eq:thetaoriginal}, and
\[
 A_-(x)\equiv A_+(x)\equiv\Theta(x)\pmod4.
\]
In particular, for either sign and every $n>0$,
\[
 [x^n]A_\pm\equiv
 \begin{cases}
  2\pmod4,&n=j^2\text{ for some }j\in\N,\\
  0\pmod4,&\text{otherwise}.
 \end{cases}
\]
\end{theorem}

\begin{proof}
Work over $\Z/4\Z$ and put $S=\sum_{j\ge1}x^{j^2}$, so $\Theta=1+2S$.
The series $B=x\Theta$ has zero constant coefficient and satisfies
$B\equiv x\pmod2$. Lemma~\ref{lem:double} gives $2S(B)=2S(x)$.
Therefore
\[
 \Theta(x\Theta(x))=1+2S(x\Theta(x))=1+2S(x)=\Theta(x).
\]
The uniqueness just proved identifies $\Plus(\Theta)=\Theta$ over this
ring. Also
\[
 \Theta^2=(1+2S)^2=1+4S+4S^2=1,
\]
so $\Theta^{-1}=\Theta$. It follows that
\[
 \Minus(\Theta)=\Plus(\Theta^{-1})^{-1}
 =\Plus(\Theta)^{-1}=\Theta^{-1}=\Theta.
\]
Compatibility with coefficient reduction identifies these series with
the reductions of the canonical integer solutions. Finally, the
coefficient of $\Theta$ at a positive index is two exactly when the
index is a square, and zero otherwise.
\end{proof}

This proves the coefficient conjecture for the selected source, with the
other substitution sign supplied by the same general construction
\cite{oeis580,oeis581}. Each equivalence is accompanied
by the complementary integer statement $4\mid a_n$ at nonsquares.
The constant term is one in both sequences; excluding $n=0$ is essential.
The modular argument only uses the form $1+2S$ and $S(0)=0$, so its
mechanism is more general than square support. The shared mechanism is included to make this manuscript self-contained.


\section{Separate Lean proof and checked scope}

The entry file is \nolinkurl{lean/Entries/A378580.lean}. Its explicit
theorem statements and proof terms use the namespace
\nolinkurl{FinalCollection.A378580}. The complete imported implementations
and their shared mathematical dependencies are included in the collection.
The following local theorem names give the proof map:

\begin{multicols}{2}\raggedright\begin{itemize}\setlength{\itemsep}{2pt}\setlength{\parsep}{0pt}
\item \nolinkurl{theta_constant_term}.
\item \nolinkurl{theta_coefficient_formula}.
\item \nolinkurl{minus_constant}.
\item \nolinkurl{quotient_identity}.
\item \nolinkurl{minus_source}.
\item \nolinkurl{minus_exists_unique}.
\item \nolinkurl{minus_unique}.
\item \nolinkurl{minus_series_mod_four}.
\item \nolinkurl{minus_congruence}.
\item \nolinkurl{minus_square_iff}.
\item \nolinkurl{minus_four_dvd_iff}.
\end{itemize}\end{multicols}

The main implementation module is \nolinkurl{Round3.ThetaSubstitution}.
The project pins Lean 4.19.0 and Mathlib
\nolinkurl{c44e0c8ee63ca166450922a373c7409c5d26b00b}.
The supplied collection report records fresh source compilation and
transitive axiom audits. All coefficient conclusions are unbounded.

\paragraph{Limits and status.}
The multiplicative unit inverse and compositional inverse are distinct. Analytic convergence and real special values are not proved.
No global novelty, exclusive priority, peer review, or OEIS acceptance is asserted.

\begin{thebibliography}{9}
\raggedright
\bibitem{oeis580}
Paul D. Hanna, \emph{OEIS A378580}, 8 January 2025.
\url{https://oeis.org/A378580}. Retrieved 24 October 2026.
\bibitem{oeis581}
Paul D. Hanna, \emph{OEIS A378581}, 8 January 2025.
\url{https://oeis.org/A378581}. Retrieved 24 October 2026.
\end{thebibliography}
\end{document}
